\section{高阶 DPM-Solver}

\subsection{二阶近似：DPM-Solver-2}

一阶近似假设 $\boldsymbol{\epsilon}_\theta$ 在整个区间内为常数。自然的改进是使用\textbf{一阶泰勒展开}（线性近似），利用两个点的噪声预测值来拟合 $\boldsymbol{\epsilon}_\theta$ 的变化趋势。

DPM-Solver-2 的算法步骤如下：

\textbf{步骤 1：计算中间点}

给定当前时间步 $s$ 和目标时间步 $t$，选取中间时间点 $r$（通常取区间中点），用一阶 DPM-Solver 计算中间值：
$$\mathbf{u} = \frac{\alpha_r}{\alpha_s} \mathbf{x}_s - \sigma_r (e^{h_1} - 1) \boldsymbol{\epsilon}_\theta(\mathbf{x}_s, s)$$

其中 $h_1 = \lambda_r - \lambda_s$。

\textbf{步骤 2：利用中间点的噪声预测更新}

在中间点 $r$ 处重新评估噪声预测 $\boldsymbol{\epsilon}_\theta(\mathbf{u}, r)$，然后用两点信息做线性插值近似积分：
$$\mathbf{x}_t = \frac{\alpha_t}{\alpha_s} \mathbf{x}_s - \sigma_t (e^{h} - 1) \boldsymbol{\epsilon}_\theta(\mathbf{x}_s, s) - \frac{\sigma_t}{2h_1}(e^{h} - 1 - h)\left[\boldsymbol{\epsilon}_\theta(\mathbf{u}, r) - \boldsymbol{\epsilon}_\theta(\mathbf{x}_s, s)\right]$$

其中 $h = \lambda_t - \lambda_s$。

\begin{importantbox}{二阶近似的关键思想}
DPM-Solver-2 每步需要 2 次网络评估（$s$ 处和 $r$ 处），但获得了二阶精度。与通用的 Heun 方法不同，DPM-Solver-2 利用了 PF-ODE 的半线性结构，因此在相同 NFE 下精度更高。
\end{importantbox}

\subsection{三阶及更高阶近似}

遵循同样的思路，可以继续做二阶泰勒展开得到 DPM-Solver-3：在区间内选取两个中间点，利用三个噪声预测值来构造二次多项式近似 $\boldsymbol{\epsilon}_\theta$ 在 $\lambda$ 空间中的变化。

\begin{knowledgebox}{递归结构}
DPM-Solver 的阶次提升具有递归结构：
\begin{itemize}
    \item \textbf{$k$ 阶 DPM-Solver} 需要 $k$ 次网络评估
    \item 先用低阶 DPM-Solver 计算 $k-1$ 个中间点
    \item 利用所有 $k$ 个噪声预测值构造 $(k-1)$ 次多项式近似
    \item 将多项式代入积分公式解析积分
\end{itemize}
实践中，2 阶和 3 阶已经足够好。超过 3 阶后，收益递减，因为噪声预测网络本身的近似误差成为瓶颈。
\end{knowledgebox}

\subsection{各阶 DPM-Solver 对比}

\begin{table}[H]
\centering
\begin{tabular}{lccc}
\toprule
\textbf{方法} & \textbf{每步 NFE} & \textbf{精度阶数} & \textbf{备注} \\
\midrule
DDPM & 1 & 约 0.5 阶 & 朴素欧拉-丸山离散化 \\
DDIM / DPM-Solver-1 & 1 & 1 阶 & 指数积分 + 零阶噪声近似 \\
DPM-Solver-2 & 2 & 2 阶 & 指数积分 + 一阶噪声近似 \\
DPM-Solver-3 & 3 & 3 阶 & 指数积分 + 二阶噪声近似 \\
\bottomrule
\end{tabular}
\caption{DPM-Solver 各阶方法的特性比较}
\end{table}

\subsection{本章小结}

高阶 DPM-Solver 通过在 $\lambda$ 空间中对噪声预测做高阶多项式近似，在每步仅增加少量网络评估的代价下获得更高精度。这种递归构造使得从一阶到三阶的提升在概念上非常自然。


